% answers.tex —— 1977 年·江西卷 · 参考答案与详细解析片段
% 由 tools/generate_projects.py 依来源总账自 1977/1977_数学真题_江西_试题.md + 1977/1977_数学真题_江西_解析正文.md + 1977/1977_数学真题_江西_答案要点.md 生成，勿手改（改源 MD 后重跑）


\begin{gktip}
固定源是题面、解答与〔评分标准〕连排的题解排印本（书页 40—49）。下文逐题独立推导，并与源解答复算对照：第一大题第（7）—（15）小项源卷括号内印有答案，与本文结论一致；第三题题面数列第三项前的短横状墨迹按源解答口径判读为 b/a（见试题源注）；第四题源解答的排列记法排印较乱，本文按通行记法重写，数值与源一致；第五题原卷排印「4=2B」按字形判读为 A=2B（见试题源注）；第六题源解答注明「图自画」，本文按题设直接推算，未配图。
\end{gktip}

\begin{answer}
\gksolhead{1}{}

\textbf{题目}（20 分）：

下列各题将答案填在括号内：

\begin{gktip}
\textbf{题号说明}：源卷第一题的 15 个小项印作「1.」—「15.」，本库规范要求小问用全角括号编号，故重排为（1）—（15），未增删题干内容。
\end{gktip}

（1） $f(x)=3x^3-5x^2+7x-5$，则 $f(0)=$\mbox{（　　）}；

（2） $\operatorname{arctg}1=$\mbox{（　　）}；

（3） $\log_a a^2=$\mbox{（　　）}；

（4） 球的面积是大圆面积的\mbox{（　　）}倍；

（5） 若 $a<0$，则 $\sqrt{a^2}\div a=$\mbox{（　　）}；

（6） $\dfrac12+\dfrac{\sqrt3}{2}i$ 的模是\mbox{（　　）}，幅角是\mbox{（　　）}，三角函数式是\mbox{（　　）}；

（7） 若 $\alpha$ 为锐角，又 $\sin 2\alpha=1$，则 $\sin\alpha=$\mbox{（　　）}；

（8） 若 $a$ 大于 $b$、$c$ 小于零，则 $ac$\mbox{（　　）}于 $bc$；

（9） 三角形的重心是三条\mbox{（　　）}线的交点；

（10） 若 $C_n^3=C_n^5$，则 $n=$\mbox{（　　）}；

（11） 直线的方程为 $y=-3x+2$，则它的斜率是\mbox{（　　）}；

（12） 正圆锥底半径长 $3$ 公分，母线长 $5$ 公分，则高是\mbox{（　　）}公分；

（13） $y=2\sin\left(3x+\dfrac{\pi}{6}\right)$ 的振幅是\mbox{（　　）}，周期是\mbox{（　　）}；

（14） $\dfrac{\sqrt3+\sqrt2}{\sqrt3-\sqrt2}$ 与 $\dfrac{\sqrt3-\sqrt2}{\sqrt3+\sqrt2}$ 的等差中项是\mbox{（　　）}，等比中项是\mbox{（　　）}；

（15） 圆的半径为 $R$，内接正方形的边长为\mbox{（　　）}，内接正三角形边长为\mbox{（　　）}。

\begin{gktip}
\textbf{源注}：原卷「圆」字排作「园」（如「正园锥」「园的半径」），本文统一转写为「圆」。原卷第（7）—（15）小项括号内排有答案，本卷一律留空。
\end{gktip}

\Needspace{4\baselineskip}
\textbf{解析}

\gksubpart{（1）}

代入 $x=0$：$f(0)=3\cdot 0^3-5\cdot 0^2+7\cdot 0-5=-5$。

\gksubpart{（2）}

$\operatorname{arctg}1$ 即 $\arctan 1$，值为 $\dfrac{\pi}{4}$。

\gksubpart{（3）}

$\log_a a^2=2$（$a>0$ 且 $a\ne1$）。

\gksubpart{（4）}

球半径 $R$ 时球面积为 $4\pi R^2$，大圆面积为 $\pi R^2$，故为大圆面积的 $4$ 倍。

\gksubpart{（5）}

\Needspace{5\baselineskip}
$a<0$ 时 $\sqrt{a^2}=|a|=-a$，故

\nopagebreak[4]
\[
\sqrt{a^2}\div a=\frac{-a}{a}=-1.
\]

\gksubpart{（6）}

模为 $\sqrt{\left(\dfrac12\right)^2+\left(\dfrac{\sqrt3}{2}\right)^2}=1$；幅角 $\varphi$ 满足 $\cos\varphi=\dfrac12$、$\sin\varphi=\dfrac{\sqrt3}{2}$，主值为 $\dfrac{\pi}{3}$；三角函数式为 $\cos\dfrac{\pi}{3}+i\sin\dfrac{\pi}{3}$。

\gksubpart{（7）}

$\sin 2\alpha=1$ 且 $\alpha$ 为锐角，故 $2\alpha=90^\circ$，$\alpha=45^\circ$，$\sin\alpha=\dfrac{\sqrt2}{2}$。

\gksubpart{（8）}

$c<0$ 时 $ac-bc=(a-b)c$，其中 $a-b>0$，故 $(a-b)c<0$，即 $ac<bc$，括号内填「小」。

\gksubpart{（9）}

三角形的重心是三条中线的交点，括号内填「中」。

\gksubpart{（10）}

由组合数性质，$C_n^3=C_n^5$（从 $n$ 个中取 $3$ 个与取 $5$ 个的组合数相等，两数不同）当且仅当 $n=3+5$，故 $n=8$。

\gksubpart{（11）}

$y=kx+b$ 的斜率为 $k$，故 $y=-3x+2$ 的斜率为 $-3$。

\gksubpart{（12）}

\Needspace{5\baselineskip}
母线 $l=5$、底半径 $r=3$，由勾股定理

\nopagebreak[4]
\[
h=\sqrt{l^2-r^2}=\sqrt{25-9}=4
\]

（公分）。

\gksubpart{（13）}

振幅 $A=2$；周期 $T=\dfrac{2\pi}{3}$（$\omega=3$）。

\gksubpart{（14）}

\Needspace{5\baselineskip}
先分母有理化：

\nopagebreak[4]
\[
\begin{aligned}
\frac{\sqrt3+\sqrt2}{\sqrt3-\sqrt2}
&=\frac{(\sqrt3+\sqrt2)^2}{(\sqrt3)^2-(\sqrt2)^2}\\
&=5+2\sqrt6,
\end{aligned}
\]

同理 $\dfrac{\sqrt3-\sqrt2}{\sqrt3+\sqrt2}=5-2\sqrt6$。两数之和为 $10$，等差中项为 $5$；两数之积为 $25-24=1$，等比中项为 $\pm1$。

\gksubpart{（15）}

内接正方形的对角线即圆的直径 $2R$，边长为 $\sqrt2\,R$；内接正三角形的边所对圆心角为 $120^\circ$，边长为 $2R\sin 60^\circ=\sqrt3\,R$。

\gkanswer{答案（1）：$-5$。}

\gkanswer{答案（2）：$\dfrac{\pi}{4}$。}

\gkanswer{答案（3）：$2$。}

\gkanswer{答案（4）：$4$。}

\gkanswer{答案（5）：$-1$。}

\gkanswer{答案（6）：模 $1$，幅角 $\dfrac{\pi}{3}$，三角函数式 $\cos\dfrac{\pi}{3}+i\sin\dfrac{\pi}{3}$。}

\gkanswer{答案（7）：$\dfrac{\sqrt2}{2}$。}

\gkanswer{答案（8）：小（$ac<bc$）。}

\gkanswer{答案（9）：中。}

\gkanswer{答案（10）：$8$。}

\gkanswer{答案（11）：$-3$。}

\gkanswer{答案（12）：$4$ 公分。}

\gkanswer{答案（13）：振幅 $2$，周期 $\dfrac{2\pi}{3}$。}

\gkanswer{答案（14）：等差中项 $5$，等比中项 $\pm1$。}

\gkanswer{答案（15）：正方形边长 $\sqrt2\,R$，正三角形边长 $\sqrt3\,R$。}

\end{answer}
\begin{answer}
\gksolhead{2}{}

\textbf{题目}（15 分）：

若方程 $5x^2-10x\cos\alpha+7\cos\alpha+6=0$ 的两根相等，试求两邻边之和为 $6$，且夹角为 $\alpha$ 的平行四边形的最大面积。

\Needspace{4\baselineskip}
\textbf{解析}

\Needspace{5\baselineskip}
两根相等，故判别式为零：

\nopagebreak[4]
\[
\begin{aligned}
\Delta&=(10\cos\alpha)^2-4\cdot5\cdot(7\cos\alpha+6)\\
&=20\left(5\cos^2\alpha-7\cos\alpha-6\right)=0,
\end{aligned}
\]

\Needspace{5\baselineskip}
即 $5\cos^2\alpha-7\cos\alpha-6=0$，分解得

\nopagebreak[4]
\[
(5\cos\alpha+3)(\cos\alpha-2)=0.
\]

\Needspace{5\baselineskip}
$\cos\alpha=2$ 舍去，故 $\cos\alpha=-\dfrac35$（此时 $\alpha$ 为钝角），

\nopagebreak[4]
\[
\sin\alpha=\sqrt{1-\cos^2\alpha}=\frac45.
\]

\Needspace{5\baselineskip}
设平行四边形一边为 $x$，则另一边为 $6-x$，夹角为 $\alpha$，面积

\nopagebreak[4]
\[
\begin{aligned}
S&=x(6-x)\sin\alpha=\frac45x(6-x)\\
&=\frac45\left[9-(3-x)^2\right]\le\frac45\cdot 9=\frac{36}{5},
\end{aligned}
\]

当 $x=3$ 时取等号。故最大面积为 $\dfrac{36}{5}$（面积单位），此时两邻边均为 $3$。

\gkanswer{答案：最大面积为 $\dfrac{36}{5}$（面积单位），此时两邻边均为 $3$。}

\end{answer}
\begin{answer}
\gksolhead{3}{}

\textbf{题目}（15 分）：

\Needspace{5\baselineskip}
若

\nopagebreak[4]
\[
5\left|\frac{a+2b-3}{b+1}\right|+3(a+3b)^2=0,
\]

求无穷数列 $ab$，$b$，$\dfrac{b}{a}$，$\cdots\cdots$ 的和。

\begin{gktip}
\textbf{源注}：源卷该数列第三项 $\dfrac{b}{a}$ 前印有一短横状墨迹，形似负号；按源解答所列数列前三项 $-27$、$-3$、$-\dfrac13$（公比 $\dfrac19$）判读应为 $\dfrac{b}{a}$，本文按此转录并留注。
\end{gktip}

\Needspace{4\baselineskip}
\textbf{解析}

\Needspace{5\baselineskip}
绝对值与平方均非负，其和为零则两者同为零：

\nopagebreak[4]
\[
a+2b-3=0,
\qquad a+3b=0.
\]

由后式 $a=-3b$，代入前式得 $-b-3=0$，故 $b=-3$，$a=9$。

\Needspace{5\baselineskip}
于是无穷数列为

\nopagebreak[4]
\[
ab,\quad b,\quad \frac{b}{a},\quad \cdots
\]

\Needspace{5\baselineskip}
具体即

\nopagebreak[4]
\[
-27,\quad -3,\quad -\frac13,\quad \cdots
\]

\Needspace{5\baselineskip}
它是等比数列，公比 $q=\dfrac{-3}{-27}=\dfrac19$，$|q|<1$，由无穷等比数列求和公式，

\nopagebreak[4]
\[
s=\frac{-27}{1-\dfrac19}=-27\times\frac98=-\frac{243}{8}=-30\frac38.
\]

\gkanswer{答案：$-\dfrac{243}{8}=-30\dfrac{3}{8}$。}

\end{answer}
\begin{answer}
\gksolhead{4}{}

\textbf{题目}（20 分）：

用 $0$、$2$、$5$、$7$、$9$ 诸数字组成没有重复数字的四位数中

（1） 大于 $5000$ 的有多少个？

（2） 是 $4$ 的倍数的有多少个？

\Needspace{4\baselineskip}
\textbf{解析}

\gksubpart{（1）}

\Needspace{5\baselineskip}
首位不能为 $0$，也不能为 $2$（否则不大于 $5000$），只能取 $5$、$7$、$9$ 之一；其余四个数位由其余四个数字全排列。故

\nopagebreak[4]
\[
3\times A_4^4=3\times24=72
\]

个。

\gksubpart{（2）}

四位数为 $4$ 的倍数当且仅当末两位数是 $4$ 的倍数。用这些数字组成的无重复数字的两位结尾只有 $20$、$52$、$72$、$92$ 四种。

\Needspace{5\baselineskip}
末尾为 $20$：前两位从 $5$、$7$、$9$ 中取两个作排列，有 $A_3^2=6$ 个；末尾为 $52$、$72$、$92$：前两位从含 $0$ 的其余三个数字中取两个作排列，首位不能为 $0$，各有 $A_3^2-A_2^1=6-2=4$ 个。故

\nopagebreak[4]
\[
A_3^2+3\left(A_3^2-A_2^1\right)=6+3\times4=18
\]

个。

\gkanswer{答案（1）：$72$ 个。}

\gkanswer{答案（2）：$18$ 个。}

\end{answer}
\begin{answer}
\gksolhead{5}{}

\textbf{题目}（15 分）：

三角形的三边为 $a$，$b$，$c$，它们的对角分别为 $A$、$B$、$C$，若 $a^2=b(b+c)$，求证 $A=2B$。

\begin{gktip}
\textbf{源注}：原卷此行铅字排印作「b(b+C)」「4=2B」——系该排印本字形与大小写混排所致（斜体 A 形近「4」、边 c 排作大写 C），按题干「对角分别为 A、B、C」的对应关系分别转写为 b(b+c) 与 A=2B。
\end{gktip}

\Needspace{4\baselineskip}
\textbf{解析}

\Needspace{5\baselineskip}
\textbf{法一（相似三角形）} 延长 $CA$ 至 $D$，使 $AD=c$，则 $CD=b+c$。因 $AB=AD=c$，$\triangle ABD$ 等腰，故

\nopagebreak[4]
\[
\angle ABD=\angle ADB.
\]

\Needspace{5\baselineskip}
由 $a^2=b(b+c)$ 得 $BC^2=AC\cdot CD$，即 $BC:CD=AC:BC$；又 $\angle ACB=\angle BCD$（公共角），故

\nopagebreak[4]
\[
\triangle ABC\sim\triangle BDC,
\]

\Needspace{5\baselineskip}
从而 $\angle ABC=\angle BDC$。$D$、$A$、$C$ 共线，$\angle BDC=\angle BDA$，于是 $\angle ABC=\angle ADB$。而 $\angle CAB$ 是 $\triangle ABD$ 在顶点 $A$ 处的外角，

\nopagebreak[4]
\[
\angle CAB=\angle ABD+\angle ADB=2\angle ADB=2\angle ABC,
\]

即 $A=2B$。

\Needspace{5\baselineskip}
\textbf{法二（余弦定理）} 由 $a^2=b^2+bc$，

\nopagebreak[4]
\[
\begin{aligned}
\cos A&=\frac{b^2+c^2-a^2}{2bc}
=\frac{c-b}{2b},\\
\cos B&=\frac{a^2+c^2-b^2}{2ac}
=\frac{b+c}{2a}.
\end{aligned}
\]

\Needspace{5\baselineskip}
于是

\nopagebreak[4]
\[
\begin{aligned}
\cos 2B&=2\cos^2B-1
=\frac{(b+c)^2}{2a^2}-1\\
&=\frac{(b+c)^2-2b(b+c)}{2b(b+c)}
=\frac{c-b}{2b}=\cos A.
\end{aligned}
\]

故 $\cos A=\cos 2B$。若 $2B=360^\circ-A$，则 $A+B=360^\circ-B>180^\circ$（因 $B<180^\circ$），与三角形内角和矛盾，故只能 $A=2B$。

\gkanswer{答案：见解析版}

\end{answer}
\begin{answer}
\gksolhead{6}{}

\textbf{题目}（15 分）：

设等腰 $\triangle ABC$ 的底边 $BC$ 在平面 $P$ 上，它的底角为 $75^\circ$，顶点 $A$ 到平面 $P$ 的距离 $AD=4$ 公分，$\triangle ABC$ 所在平面与平面 $P$ 成 $60^\circ$ 的角，求棱锥 $A-BCD$ 的体积。

\begin{gktip}
\Needspace{0.25\textheight}
\textbf{附加题}
\end{gktip}

\textbf{附加题 7．（20 分）} 已知抛物线 $y^2=x$ 上两点 $A(1,1)$ 和 $B(4,-2)$，求过 $A$、$B$ 的直线 $AB$ 与抛物线 $y^2=x$ 所围成图形的面积。

\Needspace{4\baselineskip}
\textbf{解析}

设 $D$ 为 $A$ 在平面 $P$ 上的射影，$AD=4$ 公分；作 $AE\perp BC$，垂足为 $E$。等腰三角形底边上的高同时是中线，故 $E$ 为 $BC$ 中点，$BE=EC$。

$AD\perp$ 平面 $P$，$AE$ 是斜线，由三垂线定理得 $DE\perp BC$，因此 $\angle AED$ 就是 $\triangle ABC$ 所在平面与平面 $P$ 所成二面角的平面角，为 $60^\circ$。

\Needspace{5\baselineskip}
直角三角形 $ADE$ 中：

\nopagebreak[4]
\[
ED=\frac{AD}{\tan 60^\circ}=\frac{4}{\sqrt3},
\qquad
AE=\frac{AD}{\sin 60^\circ}=\frac{8}{\sqrt3}.
\]

\Needspace{5\baselineskip}
底角为 $75^\circ$，故直角三角形 $ABE$ 中 $\angle BAE=15^\circ$，

\nopagebreak[4]
\[
BE=EC=AE\tan 15^\circ=\frac{8}{\sqrt3}(2-\sqrt3).
\]

\Needspace{5\baselineskip}
于是

\nopagebreak[4]
\[
\begin{aligned}
S_{\triangle BCD}
&=\frac12\cdot BC\cdot ED=EC\cdot ED\\
&=\frac{8}{\sqrt3}(2-\sqrt3)\times\frac{4}{\sqrt3}
=\frac{32}{3}(2-\sqrt3),
\end{aligned}
\]

\nopagebreak[4]
\[
\begin{aligned}
V&=\frac13\,S_{\triangle BCD}\cdot AD\\
&=\frac13\times\frac{32}{3}(2-\sqrt3)\times4
=\frac{128}{9}(2-\sqrt3)
\end{aligned}
\]

（立方公分）。

\textbf{附加题 7}

\Needspace{5\baselineskip}
直线 $AB$ 的斜率 $k=\dfrac{-2-1}{4-1}=-1$，方程为 $y=-x+2$，即 $x=2-y$。它与抛物线 $x=y^2$ 交于 $A(1,1)$、$B(4,-2)$，交点的纵坐标分别为 $1$ 与 $-2$。所围图形介于两交点之间，右侧以直线为界、左侧以抛物线为界：

\nopagebreak[4]
\[
\begin{aligned}
S&=\int_{-2}^{1}\left[(2-y)-y^2\right]dy
=\left[2y-\frac{y^2}{2}-\frac{y^3}{3}\right]_{-2}^{1}\\
&=\left(2-\frac12-\frac13\right)-\left(-4-2+\frac83\right)
=\frac76+\frac{10}{3}=\frac92=4\frac12.
\end{aligned}
\]

所求面积为 $\dfrac92$（面积单位），与源卷两种解法（积分法、分割求和法）的结果一致。

\gkanswer{答案：$V=\dfrac{128}{9}(2-\sqrt3)$ 立方公分。}

\end{answer}